\section{A certified cancellation in the negation constructor}
\label{sec:negation-peephole}

This section concerns a register-program identity, not a weakening of the
proof calculus.  The idea was identified in the present grammar search;
the arithmetic argument and the AST-level audit below were reconstructed
independently.  Let
\[
 \pi(x,y)=\frac{(x+y)(x+y+1)}2+y,
 \qquad \operatorname{Neg}(x)=4\pi(x,1).
\]
An adjacent prior route is recorded in the repository's
\texttt{intake-159-source-56.md} and \texttt{intake-158-159.md}:
the pairing-shift identity
$\pi(x,1)=\pi(0,x)+2$ gives
$\operatorname{Neg}(x)=4\pi(0,x)+8$.
That submission changes the arithmetic operand arrangement.  The present
identity instead retains the existing operand in scratch storage and
cancels a push immediately followed by its inverse; it still computes
$4\pi(x,1)$.  The independent whole-AX occurrence and scratch-precondition
audit below certifies this operand-retention transformation.  It does not
repair or incorporate the separate missing source-correspondence evidence
identified for the prior 255-state pairing-shift candidate.  These are
related constructor-optimization ideas with different proof obligations.

Use $t=r_1$ for the current formula, $s=r_2$ for the logical stack, and
$u=r_6,v=r_7$ for the two scratch registers.  All register values are arbitrary
natural numbers.  The exact fine tokens in the archived source give the
following transformations; omitted registers are unchanged:
\[
\begin{array}{c|c}
\text{token}&(t,s,u,v)\text{ after execution}\\\hline
L&(0,s,u+t,v)\\
M&(t+1,s,u,v)\\
N&(t+s,0,u,v)\\
P&(0,s,0,v+\pi(u,t))\\
Q&(t,s+v,u,0)\\
W&(t+4v,s,u,0).
\end{array}
\]
For $V$, the input condition $u=v=0$ gives
\[
 V(t,s,0,0)=(t,0,\pi_0(s),\pi_1(s)),
 \qquad \pi(\pi_0(s),\pi_1(s))=s.
\]
These are statements about the actual closed AST blocks, including their
break/continue levels.  They are not replacements of the source by a
higher-level pairing instruction.

\paragraph{The pair loop.}
The first inner loop of $P$ consumes $t$ into $u$ and $v$, giving
$u+t$ in scratch register~6 and adding $t$ to register~7.  If $n=u+t>0$,
the next test consumes one unit, adds one unit to~7, moves the remaining
$n-1$ units into register~1, and repeats.  The quantity $n$ decreases by
one on each complete outer round, and every inner transfer terminates.
Thus the extra accumulation is
$\sum_{j=1}^{u+t}j=(u+t)(u+t+1)/2$, in addition to the first $t$ units.
The result is exactly $v+\pi(u,t)$, with registers~1 and~6 zero.

\paragraph{The unpair loop.}
With scratch registers initially zero, one outer step of $V$ consumes one
unit of~$s$ and updates $(u,v)$ by
\[
 (u,v)\longmapsto
 \begin{cases}(u-1,v+1),&u>0,\\(v+1,0),&u=0.
 \end{cases}
\]
Both updates increase $\pi(u,v)$ by one.  Starting at $(0,0)$ and performing
$s$ such steps therefore returns the unique inverse pair of $s$.  Its outer
counter decreases at every iteration, and the inner move loop terminates.
This also proves the exact orientation of the two outputs.

\paragraph{The remaining symbolic operations.}
The whole-word replay also uses the tokens $R,S,X$ to copy an available
parameter before restoring its source.  For $c=R,S,X$, put respectively
$j=4,5,3$.  Their exact operation, without a cleanliness assumption, is
\[
 r_1'=r_1+r_j,\qquad r_6'=r_6+r_j,\qquad r_j'=0,
\]
with every other register unchanged.  The actual fine AST tests and
decrements $r_j$, increments registers~1 and~6 once on its positive branch,
and returns to that same guard.  After $k$ positive branches, the invariant
is $(r_j,r_1,r_6)=(r_j^0-k,r_1^0+k,r_6^0+k)$.  Hence the loop terminates
after exactly $r_j^0$ positive branches and has the claimed output.  The
tokens $K,U,J$ then move scratch~6 back to register $4,5,3$, respectively.
The other moves used by the replay are
\[
\begin{array}{c|rrrrrrrr}
\text{token}&L&N&O&K&U&Q&I&J\\\hline
\text{source}&1&2&6&6&6&7&7&6\\
\text{destination}&6&1&1&4&5&2&8&3.
\end{array}
\]
Each move adds its source to its destination and clears its source; its
single decreasing guard proves the statement for arbitrary destination
contents.  Finally, $H$ is exactly the same unpair walk as $V$ with its
input counter changed from register~2 to register~8.  Thus, when $r_6=r_7=0$,
\[
 r_8'=0,\qquad r_6'=\pi_0(r_8),\qquad r_7'=\pi_1(r_8),
\]
and all other registers are unchanged.  The preceding unpair invariant
proves this operation for every proof-list code, including zero.  These
facts account for every token interpreted by the symbolic AX body.

\paragraph{Cancellation theorem.}
For every $t,s$ and every assignment of the other five registers,
\[
 \boxed{\quad LNPQMVQPW\;\equiv\;LMPW
       \quad\text{when }r_6=r_7=0.\quad}
\]
Both sides terminate, preserve the stack $s$ and the other five registers,
return both scratch registers to zero, and replace the current formula by
$4\pi(t,1)$.  Indeed,
\[
\begin{array}{c|c}
\text{prefix}&(t,s,u,v)\\\hline
\text{start}&(t,s,0,0)\\
LNPQ&(0,\pi(t,s),0,0)\\
LNPQM&(1,\pi(t,s),0,0)\\
LNPQMVQ&(1,s,t,0)\\
LNPQMVQPW&(4\pi(t,1),s,0,0),
\end{array}
\]
whereas $LM$ reaches $(1,s,t,0)$ directly.  The old constructor writes a
stack cell and immediately unpairs it; the new constructor retains the
needed operand in its scratch register.
In fact the local equivalence only needs $v=0$: if $u$ is arbitrary, both
programs return $(4\pi(t+u,1),s,0,0)$.  Requiring $u=0$ gives the intended
negation of the original current formula and is already supplied at every
validated constructor site.

\paragraph{Why the precondition is part of the theorem.}
The equivalence is not asserted for arbitrary dirty scratch registers.
The audit includes unequal-output controls with nonzero scratch values.
The appropriate application is a formula-construction boundary, whose
contract already includes zero scratch registers.  Arbitrary rewrites in
unrelated control code would require separate justification.

\paragraph{Whole-AX premise audit.}
The independent checker reads the actual fine word for the routed source,
then interprets all its constructor tokens over an interned symbolic
expression DAG.  Its variables stand for arbitrary initial registers;
only registers~6 and~7 are set to zero.  A pair expression is split into its
exact two children; all other inverse-pair expressions are retained as
symbolic projections.  The rules above, ordinary register transfers, scalar
copies and the corresponding proof-stream unpair operation give the
symbolic semantics.  Each token's rule is justified by its loop invariant.
Consequently structural induction on the word makes this an all-values
source certificate, not a run on sample formula codes.

At every literal occurrence of the old nine-token negation word, the
checker requires both incoming scratch expressions to be literally zero.
It verifies the local replacement's complete nine-register output, performs
all replacements, and compares the complete output again.  The unchanged
five-token gate/cleanup suffix sees exactly the same register state.
The checker also audits the changed-formula controls using the same rule.
The JSON record lists every source offset, source/output hash and exact
number of replacements.  A separate interpreter executes the extracted
binary AST directly, with its actual decrement, break and continue
instructions, as a finite regression check of the stated token contracts.

\paragraph{Where the AX-entry invariant is established.}
The precise original-source interface is in the vendored
\texttt{L2Root254.v}, at the ProveIt commit
\texttt{fd1adfa19f8fe8da41444e3fa43b89a59d891067}.
Its definition \texttt{enc} is
\[
 \operatorname{enc}(s)=
 (\operatorname{np}s,\operatorname{tw}s,\operatorname{ws}s,
  0,\operatorname{d}s,\operatorname{h}s,0,0,\operatorname{pl}s).
\]
Lemma \texttt{prefix\_ok} proves that the fetch-and-read prefix returns
\[
 (n',\operatorname{tw}s,\operatorname{ws}s,
  a,\operatorname{d}s+x,\operatorname{h}s,0,0,p').
\]
The selector-handling prefixes in dispatcher lemmas \texttt{disp\_1},
\texttt{disp\_2}, and \texttt{disp\_ax} subsequently decrement only
register~3 before entering the selected command.  In the AX case
of \texttt{root\_sem}, the call to \texttt{ax1} therefore discharges both
scratch-zero premises by reflexivity.  That theorem returns
\texttt{enc} of the next enumerator state after every nonhalting command;
\texttt{root\_step\_gen}, \texttt{root\_step\_det}, and
\texttt{iter\_halts\_dag} carry this interface through command iteration.
The initial relation is \texttt{Main254.R0\_enc}.

The accompanying \texttt{ProgLib254.v} has the original corresponding
contracts \texttt{w198}, \texttt{w200}, \texttt{w193} for $R,S,X$;
\texttt{unpair115}, \texttt{unpair116} for $H,V$;
\texttt{pair162} for $P$; and \texttt{w117} for $W$.
Some fine ASTs reassociate the original source and exchange adjacent
increments.  Consequently their validity is established here by direct
inspection of their actual control graphs and the displayed invariants,
rather than by asserting literal identity with a differently grouped Coq
term.  The archived sources' Git blob IDs were compared exactly with the
objects named by the pinned current-tree commit.  The input manifest records
those identities.  Applying the command-boundary theorem to the routed
compiler still uses its existing source-equivalence certificate; no new
claim about an arbitrary internal instruction pointer is made.

\paragraph{Transfer to proof search.}
The original and optimized AX bodies emit exactly the same numerical rows
and the same numerical theory bundle.  The logical theory, parameter gate,
formula codec, proof-candidate stream, and conclusion code are therefore
unchanged.  The stated command invariant giving zero scratch at AX entry
transfers the identity through every AX invocation by induction on the
command trace.  The exact binary dispatcher and blank-start bootstrap
must additionally be verified for the newly compiled table.  The present
lemma discharges the new register-source obligation; it does not replace
those machine-layer obligations.
